\subsection{连续单射或满射线性映射的扰动}

在本节中，我们采用下列约定：若 E 是有限维向量空间，\(\dim(\mathrm{E})\) 表示它的维数；若 E 是无穷维向量空间，则令 \(\dim(\mathrm{E}) = +\infty \in \overline{\mathbf{R}}\) 。若 u 是核或余核为有限维的线性映射，则令 \(\operatorname{ind}(u) = \dim \operatorname{Coker}(u) - \dim \operatorname{Ker}(u)\) ，其中的运算在 \(\overline{R}\) 中进行。

设 E 和 F 是赋范空间。用 \(\mathcal{M}(\mathrm{E};\mathrm{F})\) 表示从 E 到 F 中的严格单形态射所成的集，用 \(\mathcal{Q}\mathcal{M}(\mathrm{E};\mathrm{F})\) 表示从 E 到 F 中的、核为有限维的严格态射所成的集。

命题 10. --- 设 E 和 F 是赋范空间。集 \(\mathcal{M}(\mathrm{E};\mathrm{F})\) 在 \(\mathcal{L}(\mathrm{E};\mathrm{F})\) 中是开的。它是 \(\mathcal{L}(\mathrm{E};\mathrm{F})\) 中单射映射所成的集的内部。

用 A 表示从 E 到 F 中的连续单射线性映射所成的集。为使映射 \(u \in A\) 是严格态射，必须且只需它的余范数 \(((u))\) 严格为正（III, p. 62, 注记）。而 \(((u))\) 是从 \(u\) 到 A 在 \(\mathcal{L}(\mathrm{E};\mathrm{F})\) 中的补集的距离（III, p. 65, 命题 9）。命题得证。

命题 11. --- 设 E 和 F 是巴拿赫空间。集 \(\mathcal{Q}\mathcal{M}(\mathrm{E};\mathrm{F})\) 在 \(\mathcal{L}(\mathrm{E};\mathrm{F})\) 中是开的。它是 \(\mathcal{L}(\mathrm{E};\mathrm{F})\) 中核为有限维的映射所成的集的内部。

用 A 表示从 E 到 F 中的、核为有限维的连续线性映射所成的集。

设 u 是 \(\mathcal{Q}\mathcal{M}(\mathrm{E};\mathrm{F})\) 的一个元素。此时有 \(((u)) > 0\) （III, p. 62, 注记）。每个元素 \(v \in \mathcal{L}(\mathrm{E};\mathrm{F})\) （与 u 的距离 \(<((u))\) ）都属于 A（III, p. 65, 命题 9），故 u 是 A 的一个内点。

反之，设 \(u\) 是 A 的一个不是严格态射的元素。我们有 \(((u)) = 0\) （III, p. 62, 注记）。设 \(v\) 是 \(\mathcal{L}(\mathrm{E};\mathrm{F})\) 的一个与 \(u\) 相差一个有限秩线性映射的元素；由 III, p. 40 的推论 1，态射 \(v\) 不可能是具有闭像的严格态射；由于从 E 到 F 中的严格态射在 F 中有闭像（EVT, IV, p. 28, 定理 1），这意味着 v 不是严格态射。于是有 \(((v)) = 0\) 。设 \(\varepsilon\) 为一实数，\textgreater{} 0。借助前面的考察，III, p. 66 的引理 3 使我们可以归纳地构造元素的一个序列 \((u_{m})_{m \in \mathbf{N}}\) （\(\mathcal{L}(\mathrm{E}; \mathrm{F})\) 中的），它满足下列条件：

\begin{enumerate}
\def\labelenumi{(\roman{enumi})}
\item
  有 \(u_{0}=u\) ；
\item
  对每个 \(m \geqslant 0\) ，\(u_{m+1} - u_m\) 是秩为 1 且范数 \(\leqslant 2^{-m-1}\varepsilon\) 的连续线性映射；
\item
  对每个 \(m \geqslant 0\) ，\(u_m\) 的核维数为 \(n + m\) ，且包含于 \(u_{m+1}\) 的核中。
\end{enumerate}

序列 \((u_{m})\) 是巴拿赫空间 \(\mathcal{L}(\mathrm{E};\mathrm{F})\) 中的柯西序列。设 \(u'\) 是它的极限。\(u'\) 的核包含 \(u_{m}\) 的核（对一切 \(m \geqslant 0\) ）；它是无穷维的。由于

\[
\left\| u ^ {\prime} - u \right\| \leqslant \sum_ {m = 0} ^ {\infty} \left\| u _ {m + 1} - u _ {m} \right\| \leqslant \varepsilon \sum_ {m = 0} ^ {\infty} 2 ^ {- m - 1} = \varepsilon ,
\]

u 到 A 的补集的距离小于 \(\varepsilon\) 。由于这对每个 \(\varepsilon > 0\) 都成立，由此得出 u 不是 A 的内点。

命题 12. --- 设 E 和 F 是巴拿赫空间，u 是 \(\mathcal{Q}\mathcal{M}(\mathrm{E};\mathrm{F})\) 的一个元素。每个 \(v \in \mathcal{L}(\mathrm{E};\mathrm{F})\) （满足 \(\|v - u\| < ((u))\) ）都属于 \(\mathcal{Q}\mathcal{M}(\mathrm{E};\mathrm{F})\) 且满足关系式

\[
\begin{array}{c} \dim \operatorname{Ker} (v) \leqslant \dim \operatorname{Ker} (u), \\ \dim \operatorname{Coker} (v) \leqslant \dim \operatorname{Coker} (u), \\ \operatorname{ind} (v) = \operatorname{ind} (u). \end{array}
\]

由于 \(u\) 是严格的，我们有 \(((u)) > 0\) 。用 B 表示以 \(v\) 为中心、\(((u))\) 为半径的开球（在 \(\mathcal{L}(\mathrm{E};\mathrm{F})\) 中）。对每个 \(v \in \mathrm{B}\) ，我们有 \(\dim \operatorname{Ker}(v) \leqslant \dim \operatorname{Ker}(u)\) （III, p. 65, 命题 9）且 \(v \in \mathcal{Q}\mathcal{M}(\mathrm{E};\mathrm{F})\) （命题 11）。

对 \(r \in \mathbf{Z} \cup \{+\infty\}\) ，用 \(B_r\) 表示由元素 \(v \in B\) （满足 \(\text{ind}(v) = r\) ）所成的集。若 \(r \in \mathbf{Z}\) ，则集 \(B_r\) 是从 E 到 F 的、指标为 \(r\) 且属于 B 的弗雷德霍姆映射所成的集（III, p. 52, 命题 11）；由 III, p. 58 的定理 1，它在 B 中是开的。

我们来证明诸集 \(B_{r}\) 在 B 中是闭的。设 \(v \in B\) 是附着于 \(B_{r}\) 的一点。由于诸集 \(B_{s}\) （\(s \in Z\) ）在 B 中是开的且两两不相交，我们有 \(v \in B_{r}\) 或 \(v \in B_{+\infty}\) 。

设 \(v \in B_{+\infty}\) 。用 n 表示 \(\operatorname{Ker}(u)\) 的维数。选取 F 的一个维数为 \(r + 2n + 1\) 的向量子空间 T，它与 \(\operatorname{Im}(v)\) 的交缩减为 0；再选取 \(\operatorname{Ker}(v)\) 在 E 中的一个拓扑补 S（III, p. 55, 命题 1）。考察连续线性映射 \(f: (s, t) \mapsto v(s) + t\) （从 S × T 到 F）。它是单射的。线性映射 v 是严格态射，因为 v 属于 B ⊂ \(\mathcal{Q}\mathcal{M}(\mathrm{E};\mathrm{F})\) 。于是 v 的像是闭的（EVT, IV, p. 28, 定理 1）。由 III, p. 39 的命题 1，v 在 S 上的限制是从 S 到 F 的具有闭像的严格态射，而 f 是从 S × T 到 F 的具有闭像的严格态射。

由命题 10，存在 \(v\) 在 B 中的一个邻域 U，使得对每个 \(w \in \mathrm{U}\) ，线性映射 \((s,t) \mapsto w(s) + t\) （从 \(\mathrm{S} \times \mathrm{T}\) 到 F 中）是单射的。设 \(w\) 是 U 的一个元素。此时有 \(\operatorname{Ker}(w) \cap \mathrm{S} = \{0\}\) ，从而 \(w\) 通过限制并过渡到商空间诱导一个从 \(\operatorname{Ker}(w)\) 到 E/S 中的单射线性映射，这蕴含 \(\operatorname{Ker}(w)\) 的维数至多为 \(n\) 。我们还有 \(w(\mathrm{S}) \cap \mathrm{T} = \{0\}\) ，由此

\[
\operatorname{codim} _ {\mathrm{F}} (\operatorname{Im} (w)) \geqslant \dim (\mathrm{T}) - \operatorname{codim} _ {\mathrm{E}} (\mathrm{S}) = r + n + 1,
\]

这蕴含 \(\operatorname{ind}(w) \geqslant r + 1\) ，且与 v 附着于 B \(_{r}\) 这一假设相矛盾。

若 \(r\) 是 \(\mathbf{Z}\) 中使 \(B_r\) 非空的元素，则我们有 \(B_r = B\) ，因为 \(B_r\) 在 \(B\) 中既开又闭且 \(B\) 连通。若 \(B_r\) 对每个 \(r \in \mathbf{Z}\) 都是空的，则有 \(\text{ind}(v) = +\infty\) （对一切 \(v \in B\) ）。于是映射 \(v \mapsto \text{ind}(v)\) 在 \(B\) 上是常数。最后，对每个 \(v \in B\) ，我们有

\[
\begin{array}{c} \dim \operatorname{Coker} (v) = \operatorname{ind} (v) + \dim \operatorname{Ker} (v) \\ \leqslant \operatorname{ind} (u) + \dim \operatorname{Ker} (u) = \dim \operatorname{Coker} (u). \end{array}
\]

证明到此结束。

推论 1. --- 由 \(u \mapsto \dim \operatorname{Ker}(u)\) 与 \(u \mapsto \dim \operatorname{Coker}(u)\) 在 \(\mathcal{Q}\mathcal{M}(\mathrm{E};\mathrm{F})\) 上定义的函数是上半连续的。函数 \(u \mapsto \operatorname{ind}(u)\) 在 \(\mathcal{Q}\mathcal{M}(\mathrm{E};\mathrm{F})\) 上是局部常数的。

推论 2. --- 函数 \(u \mapsto \dim \operatorname{Coker}(u)\) 在从 E 到 F 中的严格单形态射所成的集 \(\mathcal{M}(\mathrm{E};\mathrm{F})\) 上是局部常数的。

事实上，有 \(\dim \operatorname{Coker}(u) = \operatorname{ind}(u)\) （对 \(u \in \mathcal{M}(\mathrm{E};\mathrm{F})\) ）。

引理 4. --- 设 E 和 F 是巴拿赫空间。为使 \(\mathcal{L}(\mathrm{E};\mathrm{F})\) 的元素 u 是从 E 到 F 中的严格态射，必须且只需 \(^{t}u\) 是从 F' 到 E' 中的严格态射。此时有 \(\dim \operatorname{Coker}(^{t}u)=\dim\operatorname{Ker}(u)\) 与 \(\dim\operatorname{Ker}(^{t}u)=\dim\operatorname{Coker}(u)\) 。

为使 u 是严格态射，必须且只需 \(^{t}u\) 是（EVT, IV, p. 29, 推论 3）；此时 u 的像是闭的（EVT, IV, p. 28, 定理 1），且向量空间 \(\mathrm{Ker}(^{t}u)\) （相应地，\(\mathrm{Coker}(^{t}u)\) ）典范地同构于赋范空间 \(\mathrm{Coker}(u)\) （相应地，\(\mathrm{Ker}(u)\) ）的对偶（依照 EVT, IV, p. 27, 命题 2）。引理得证。

设 E 和 F 是巴拿赫空间。用 \(\mathcal{E}(\mathrm{E};\mathrm{F})\) 表示从 E 到 F 中的连续满射线性映射所成的集，用 \(\mathcal{Q}\mathcal{E}(\mathrm{E};\mathrm{F})\) 表示从 E 到 F 中的、像有有限余维数的连续线性映射所成的集。\(\mathcal{Q}\mathcal{E}(\mathrm{E};\mathrm{F})\) 的每个元素都是具有闭像的严格态射（III, p. 52, 引理 6）。由引理 4 得出，\(\mathcal{E}(\mathrm{E};\mathrm{F})\) 与 \(\mathcal{Q}\mathcal{E}(\mathrm{E};\mathrm{F})\) 分别是 \(\mathcal{M}(\mathrm{F}^{\prime};\mathrm{E}^{\prime})\) 与 \(\mathcal{Q}\mathcal{M}(\mathrm{F}^{\prime};\mathrm{E}^{\prime})\) 在连续映射 \(u\mapsto{}^{t}u\) （从 \(\mathcal{L}(\mathrm{E};\mathrm{F})\) 到 \(\mathcal{L}(\mathrm{F}^{\prime};\mathrm{E}^{\prime})\) 中）下的原像。

命题 13. --- 设 E 和 F 是巴拿赫空间。集 \(\mathcal{E}(\mathrm{E};\mathrm{F})\) 与 \(\mathcal{Q}\mathcal{E}(\mathrm{E};\mathrm{F})\) 在 \(\mathcal{L}(\mathrm{E};\mathrm{F})\) 中是开的。更精确地说，若 \(u\) 是 \(\mathcal{Q}\mathcal{E}(\mathrm{E};\mathrm{F})\) 的元素，而 \(v\) 是 \(\mathcal{L}(\mathrm{E};\mathrm{F})\) 中满足 \(\| v - u\| < ((u))\) 的元素，则我们有 \(v \in \mathcal{Q}\mathcal{E}(\mathrm{E};\mathrm{F})\) ，且

\(\dim \operatorname{Ker}(v) \leqslant \dim \operatorname{Ker}(u), \quad \dim \operatorname{Coker}(v) \leqslant \dim \operatorname{Coker}(u),\)

\[
\operatorname{ind} (v) = \operatorname{ind} (u).
\]

我们前面已经看到 \(^{t}u \in \mathcal{Q}\mathcal{M}(F';E')\) 。此外，对每个元素 \(v\) （在 \(\mathcal{L}(E;F)\) 中），有 \(\|^{t}v - ^{t}u\| = \|v - u\|\) 与 \(((^{t}u)) = ((u))\) （EVT, IV, p. 7, 命题 8 及 III, p. 63, 命题 8）。由命题 12，从这些关系式得出：若 \(\|v - u\| < ((u))\) ，则 \(^{t}v\) 属于 \(\mathcal{Q}\mathcal{M}(F';E')\) ，且有不等式

\[
\dim \operatorname{Ker} (^ {t} v) \leqslant \dim \operatorname{Ker} (^ {t} u), \quad \dim \operatorname{Coker} (^ {t} v) \leqslant \dim \operatorname{Coker} (^ {t} u)
\]

以及等式 \(\mathrm{ind}(^{t}v)=\mathrm{ind}(^{t}u)\) 。于是命题由引理 4 得出。

推论 1. --- 由 \(u \mapsto \dim \operatorname{Ker}(u)\) 与 \(u \mapsto \dim \operatorname{Coker}(u)\) 在 \(\mathcal{Q}\mathcal{E}(\mathrm{E};\mathrm{F})\) 上定义的函数是上半连续的。函数 \(u \mapsto \operatorname{ind}(u)\) 在 \(\mathcal{Q}\mathcal{E}(\mathrm{E};\mathrm{F})\) 上是局部常数的。

推论 2. --- 由 \(u \mapsto \dim \operatorname{Ker}(u)\) 定义的函数在 \(\mathcal{E}(\mathrm{E};\mathrm{F})\) 上是局部常数的。

事实上，有 \(\dim \operatorname{Ker}(u) = -\operatorname{ind}(u)\) （对一切 \(u \in \mathcal{E}(\mathrm{E};\mathrm{F})\) ）。

